\section*{Hoofdstuk \ref{ch:b1:p-block-13-15} --- Het p-blok I: de groepen van boor, koolstof en stikstof}

\begin{solution}{exo:b1:p-block-13-15:1}
\ce{NH4NO3}: $-$III in \ce{NH4+}, $+$V in \ce{NO3-}. \ce{N2O5}: $+$V.
\ce{HNO2}: $+$III. \ce{N2H4}: $-$II. \ce{Mg3N2}: $-$III.
\end{solution}

\begin{solution}{exo:b1:p-block-13-15:2}
\ce{BF3}: boor met drie bindingen en zes elektronen, vlak trigonaal.
\ce{NH3}: drie bindingen en een \omterm{def:b1:lewis-resonance-vsepr:lewis}{vrij paar}. Het \omterm{def:b1:lewis-resonance-vsepr:lewis}{vrije paar} van N vormt een
\omterm{def:b1:lewis-resonance-vsepr:lewis-acid}{datieve binding} met B: \ce{BF3} is het \omterm{def:b1:lewis-resonance-vsepr:lewis-acid}{lewiszuur}, \ce{NH3} de \omterm{def:b1:lewis-resonance-vsepr:lewis-acid}{lewisbase};
in het adduct zijn beide atomen tetraëdrisch.
\end{solution}

\begin{solution}{exo:b1:p-block-13-15:3}
\ce{CO2 + 2NaOH -> Na2CO3 + H2O}; \ce{SiO2 + 2NaOH -> Na2SiO3 + H2O};
\ce{P4O10 + 12NaOH -> 4Na3PO4 + 6H2O}; \ce{Al2O3 + 6HCl -> 2AlCl3 + 3H2O};
\ce{Al2O3 + 2NaOH + 3H2O -> 2NaAl(OH)4}.
\end{solution}

\begin{solution}{exo:b1:p-block-13-15:4}
$\Delta_r G^\circ = 2 \times (-16.45) = \qty{-32.9}{kJ/mol}$, $\log K =
32.9/5.708 = 5.76$, $K = 10^{5.76}$.
\end{solution}

\begin{solution}{exo:b1:p-block-13-15:5}
Ring van zes \ce{SiO3^2-}-eenheden: \ce{Si6O18^12-}. Dubbele keten, per
vier siliciumatomen: twee met twee gedeelde hoekpunten (3 O, lading $-2$
elk), twee met drie ($2.5$ O, lading $-1$ elk): \ce{Si4O11^6-}.
\end{solution}

\begin{solution}{exo:b1:p-block-13-15:6}
\qty{550}{kg} \ce{Al(OH)3} is $550/78.0 = \qty{7.05}{kmol}$, wat
\qty{3.53}{kmol} \ce{Al2O3} geeft, \qty{360}{kg}; bij \qty{92}{\percent}:
\qty{331}{kg}.
\end{solution}

\begin{solution}{exo:b1:p-block-13-15:7}
Boor heeft een lege orbitaal en neemt het \omterm{def:b1:lewis-resonance-vsepr:lewis}{vrije paar} op van een
hydroxide-ion dat aan water onttrokken wordt:
\ce{B(OH)3 + 2H2O <=> B(OH)4- + H3O+}. Het vrijgekomen proton komt van een
watermolecuul, niet van het boorzuur.
\end{solution}

\begin{solution}{exo:b1:p-block-13-15:8}
N gaat van $-$III naar $+$II, vijf elektronen per \ce{NH3}; elke \ce{O2}
neemt er vier op. $4 \times 5 = 5 \times 4$: \ce{4NH3 + 5O2 -> 4NO + 6H2O}.
\end{solution}

\begin{solution}{exo:b1:p-block-13-15:9}
\ce{NH3 + HNO3 -> NH4NO3}. Eén ton is $10^6/80.0 = \qty{1.25e4}{mol}$;
\qty{1.25e4}{mol} ammoniak om te neutraliseren en \qty{1.25e4}{mol} om het
salpeterzuur te maken: $2.50 \times 10^4 \times 17.0 = \qty{425}{kg}$.
\end{solution}

\begin{solution}{exo:b1:p-block-13-15:10}
Atomen van de tweede \omterm{def:b1:periodicity:table}{periode} zijn klein: hun \textit{p}-orbitalen
overlappen zijdelings goed en vormen sterke $\pi$-bindingen, zodat
\ce{N#N} en \ce{O=C=O} kleine, stabiele moleculen zijn. Atomen van de
derde \omterm{def:b1:periodicity:table}{periode} zijn groter; hun $\pi$-overlap is zwak, en ze verkiezen meer
enkelvoudige bindingen: \ce{P4}-tetraëders, het \ce{SiO4}-netwerk.
\end{solution}

\begin{solution}{exo:b1:p-block-13-15:11}
$E^\circ(\ce{PbO2}/\ce{Pb^2+}) = 1.46 > E^\circ(\ce{Cl2}/\ce{Cl-}) = 1.36$~V:
\ce{PbO2 + 4H+ + 2Cl- -> Pb^2+ + Cl2 + 2H2O}. Lood(IV) is instabiel ten
opzichte van lood(II) (inert paar); silicium(II) is geen stabiele toestand
en \ce{SiO2} is geen \omterm{def:b1:nernst:couple}{oxidator}.
\end{solution}

\begin{solution}{exo:b1:p-block-13-15:12}
$160 \times 17.0/14.0 = \qty{194}{Mt}$ ammoniak, met $194 \times
3.0/17.0 = \qty{34}{Mt}$ waterstof, vooral gemaakt uit aardgas (methaan en
stoom).
\end{solution}

\begin{solution}{pb:b1:p-block-13-15:1}
\textbf{1.} :N\,$\equiv$\,N: --- een drievoudige binding, zeer sterk, zonder
polariteit.
\textbf{2.} 0, $-$III, $+$II, $+$IV, $+$V, en $-$III/$+$V in \ce{NH4NO3}.
\textbf{3.} Bacteriën op de wortels van vlinderbloemigen, bliksem, het
haber-boschproces.
\textbf{4.} Het verbreken van de drievoudige binding vraagt een
\omterm{def:b1:elementary-steps:catalyst}{katalysator} die planten niet hebben; bacteriën met het enzym nitrogenase
wel.
\textbf{5.} Stikstof wordt gereduceerd (0 naar $-$III), waterstof
geoxideerd (0 naar $+$I).
\textbf{6.} \ce{N2 + 3H2 <=> 2NH3}.
\textbf{7.} $K = 10^{2 \times 16.45/5.708} = 10^{5.76}$.
\textbf{8.} Bij \qty{25}{\celsius} is de reactie veel te traag; de
\omterm{def:b1:elementary-steps:catalyst}{katalysator} werkt alleen warm, en de omstandigheden zijn een compromis
dat in het deel van jaar~2 besproken wordt.
\textbf{9.} \qty{58.8}{kmol} ammoniak: \qty{29.41}{kmol} \ce{N2},
\qty{824}{kg}; \qty{88.2}{kmol} \ce{H2}, \qty{176}{kg}.
\textbf{10.} $V(\ce{N2}) = 29.4 \times 10^3 \times 8.314 \times 298.15/10^5 =
\qty{729}{m^3}$; lucht $729/0.78 = \qty{935}{m^3}$.
\textbf{11.} Uit aardgas (stoomreforming van methaan), met \ce{CO2} als
bijproduct.
\textbf{12.} In de bodem geïnjecteerd is het een goedkope kunstmest, maar
het is een giftig gas; vaste stoffen (ureum, ammoniumnitraat) zijn
gemakkelijker op te slaan en te verspreiden.
\textbf{13.} N: $-$III $\to$ $+$II (5 elektronen), O: 0 $\to$ $-$II (4 per
\ce{O2}): \ce{4NH3 + 5O2 -> 4NO + 6H2O}.
\textbf{14.} \ce{2NO + O2 -> 2NO2}.
\textbf{15.} \ce{3NO2 + H2O -> 2HNO3 + NO}: een \omterm{def:b1:e-ph-diagrams:disproportionation}{disproportionering} (N
$+$IV naar $+$V en $+$II).
\textbf{16.} \ce{NH3 + 2O2 -> HNO3 + H2O}.
\textbf{17.} $2 \times 58.8 \times 32.0 = \qty{3.76}{t}$.
\textbf{18.} Het maakt de oxidatie snel en selectief voor NO, in plaats
van het stabielere \ce{N2}.
\textbf{19.} \ce{NH3 + HNO3 -> NH4NO3}.
\textbf{20.} De helft.
\textbf{21.} \qty{29.4}{kmol} \ce{NH4NO3}, \qty{2.35}{t}.
\textbf{22.} $28.0/80.0 = \qty{35}{\percent}$.
\textbf{23.} Het bevat een \omterm{def:b1:nernst:couple}{oxidator} (nitraat) en een \omterm{def:b1:nernst:couple}{reductor} (ammonium)
in hetzelfde \omterm{def:b1:crystals-metals:crystal}{kristal}: verhit of gemengd met brandstoffen kan het
explosief ontleden.
\textbf{24.} $58.8 \times 63.0 =$ \textbf{\qty{3.7}{t} salpeterzuur per ton
ammoniak}.
\end{solution}
